\documentclass[11pt]{article}
\usepackage[a4paper,margin=1in]{geometry}
\usepackage[T1]{fontenc}
\usepackage[utf8]{inputenc}
\usepackage{lmodern}
\usepackage[numbers,sort&compress]{natbib}
\usepackage{setspace}
\usepackage{enumitem}
\usepackage[colorlinks=true,linkcolor=blue,citecolor=blue,urlcolor=blue]{hyperref}
\setstretch{1.0}
\setlength{\parindent}{0pt}
\setlength{\parskip}{0.5em}

\begin{document}

\begin{center}
{\Large \textbf{AIAA3053 Project Proposal}}\\[0.6em]
\textbf{Course:} AIAA3053 Reinforcement Learning Principles and Methods\\
\textbf{Group:} Group 8\\
\textbf{Track:} Track A\\
\textbf{Project Topic:} Game Playing with Q-Learning\\
\textbf{Team Members:} Chi Kit WONG, Haodong XU, Sida ZHANG
\end{center}

\section*{1. Problem Statement}
This project studies \textbf{Game Playing with Q-Learning} on the \textbf{FrozenLake} environment, a discrete grid-world task in which an agent must reach a goal while avoiding terminal hole states. The problem is a suitable course project because it captures several core reinforcement learning challenges in a controlled setting: sparse rewards, sequential decision making, exploration versus exploitation, and long-term value estimation \cite{SuttonBarto2018}.

The significance of the project is twofold. First, Q-Learning is one of the most fundamental model-free reinforcement learning algorithms, so implementing it correctly is directly aligned with the course objective of understanding reinforcement learning principles \cite{WatkinsDayan1992}. Second, FrozenLake is small enough to make the learned policy interpretable, which allows us to analyze not only whether the method works but also how environment uncertainty changes the agent's behavior \cite{GymnasiumDocs}.

The research gap addressed by this project is not a new algorithmic gap in the literature, but an empirical and educational one. Many introductions to Q-Learning show only a single successful run in a simple environment. That is insufficient for understanding how stochastic transitions and hyperparameter choices affect convergence, stability, and reproducibility. Our project focuses on that gap by comparing deterministic and stochastic dynamics and by studying the effect of the learning rate across multiple random seeds.

\section*{2. Related Work}
Tabular Q-Learning is a classical reinforcement learning method and remains the standard baseline for small discrete Markov decision processes \cite{WatkinsDayan1992,SuttonBarto2018}. Existing textbook and tutorial treatments usually present the algorithm in idealized settings where the state space is small and the update rule can be inspected directly. Benchmark environments such as FrozenLake are commonly used for this purpose because they make policy quality easy to visualize and evaluation straightforward \cite{GymnasiumDocs}.

Our project differs from a basic tutorial-style implementation in three ways. First, we explicitly compare deterministic and stochastic versions of the same task instead of reporting only one environment setting. Second, we include a small but structured hyperparameter study centered on the learning rate, which is often mentioned but not analyzed carefully in introductory demonstrations. Third, we evaluate results across multiple random seeds and summarize both performance and variance, so the conclusions are based on repeatable experimental evidence rather than a single lucky run.

\section*{3. Methodology}
We will implement standard tabular \textbf{Q-Learning} with an \textbf{epsilon-greedy} exploration strategy in Python. The agent will maintain a Q-table over discrete states and actions and update it with the temporal-difference target after each transition. The implementation will include environment interaction, training and evaluation loops, configurable hyperparameters, logging, and result visualization. We plan to use \texttt{FrozenLake-v1} when available, with a compatible fallback implementation if package availability becomes a practical issue.

The implementation and experiment design will follow four steps:
\begin{enumerate}[leftmargin=1.5em]
\item build a complete training pipeline for tabular Q-Learning on FrozenLake;
\item run a deterministic baseline to verify correctness and expected learning behavior;
\item run a stochastic version with slippery dynamics to test robustness under transition noise;
\item conduct a learning-rate ablation, for example with $\alpha \in \{0.01, 0.10, 0.50\}$, to compare convergence speed and final performance.
\end{enumerate}

\section*{4. Evaluation Plan}
We will evaluate the project using both quantitative metrics and qualitative inspection of the learned policy. The main metrics will be training reward, greedy evaluation success rate, final success rate, and variance across random seeds. We will also compare approximate convergence speed by checking how quickly each setting reaches a stable success level.

The experimental settings will include deterministic FrozenLake, stochastic FrozenLake, and a learning-rate ablation in the stochastic case. Each condition will be repeated with multiple random seeds to improve reliability. We will report results using learning curves, evaluation curves, summary tables, and policy visualizations.

The expected results are that the deterministic setting will be solved reliably and quickly, while the stochastic setting will be harder and show lower success rates and higher variance. We also expect the learning-rate study to show a clear tradeoff: a very small learning rate will learn too slowly, while a better-tuned rate should improve convergence and final performance. These results should provide a concise but rigorous demonstration of how Q-Learning behaves under different environment and optimization settings.

\bibliographystyle{IEEEtran}
\bibliography{references}

\end{document}
